\section{THALAMUS: Intelligent Ingestion and Namespace Isolation for Hetero\-geneous Knowledge Graphs}

\textbf{Author}: Bryan Alexander Buchorn  
\textbf{Affiliation}: Independent Researcher, Las Vegas, NV, USA  
\textbf{Status}: v2.52.0 (Phase 172 (STRB) COMPLETE)
\textbf{Date}: May 2, 2026

---

\subsubsection{Abstract}
We present \textbf{THALAMUS}, an intelligent ingestion prepr\-ocessing pipeline designed to address the structural and semantic incon\-sistencies inherent in high-velocity, heter\-ogeneous Knowledge Graph (KG) streams. THALAMUS implements a composable archi\-tecture for entity norma\-lization, bidir\-ectional dedup\-lication, and ontology mapping. Crucially, we introduce a \textbf{Namespace Isolation} protocol that prevents "identity collapse" across data modalities (e.g., text vs. sensors) by enforcing strict prefixing and domain-specific validation. To handle the compu\-tational demands of real-time streaming, the v2.52.0 release introduces a \textbf{Parallel Ingestion Optimi\-zation} that decouples CPU-bound norma\-lization tasks from the graph's global write-lock. Benchmark results show an \textbf{850\% throughput improvement} (from 1,200 to 11,500 triples/sec) while enabling linear throughput scaling across multi-core archi\-tectures without degrading reasoning latency. As of v2.52.0, THALAMUS has been extended with a \texttt{/build} hot-reload endpoint enabling runtime CSV ingestion without server restart, and a \texttt{ResearchAgent} (Phase 51) feeds proposed edges back into the pipeline after human approval, closing the loop between autonomous hypothesis generation and structured knowledge ingestion; 1,357 tests now cover the full THALAMUS pipeline including streaming, namespace isolation, and STDP discr\-etization.

\subsubsection{1. Introd\-uction}
The "GIGO" (Garbage In, Garbage Out) principle is the primary failure mode for autonomous reasoning engines. When unrelated concepts share an ID, or when a single entity appears under multiple aliases, the graph's structural consensus (DSCF) and attention mechanisms (CSA) fail. THALAMUS acts as the "Intelligent Gatekeeper," ensuring all data is pre-aligned to a canonical repre\-sentation.

\subsubsection{2. Methodology}

\paragraph{2.1 Normal\-ization and Dedupl\-ication}
THALAMUS maintains a stateful mapping $\mathcal{M}: \{a_1, a_2, \dots, a_n\} \to e_{canonical}$. All incoming triples are projected through $\mathcal{M}$ before ingestion. This dedup\-lication process \cite{doan2012principles} utilizes both exact string matching and n-gram based fuzzy resolution.

\paragraph{2.2 Modal Isolation}
In multi-modal graphs, "Identity Collapse" occurs when a symbolic entity (text) and a physical signal (sensor) share a label. We formalize the isolation rule $\mathcal{I}$:
\begin{equation}
\mathcal{I}(id, mode) = \text{prefix}(mode) \mathbin{\|} id
\end{equation}
This ensures topological separation between disparate data layers while allowing explicit cross-modal edges [Buchorn, 2026] to bridge them.

\subsubsection{3. Scalability: Parallel Prepro\-cessing (v2.52.0)}
We define a two-stage ingestion protocol:
\begin{enumerate}
\item \textbf{Map-Stage}: $N$ worker threads normalize batches of $B$ triples in parallel. CPU-bound string work is performed outside the graph mutex.
\item \textbf{Commit-Stage}: The master thread performs a bulk-addition to the adjacency list under a single lock acquisition.
\end{enumerate}
This removes the $O(N)$ string-processing bottleneck from the critical path, unblocking query readers during high-velocity bursts.

\subsubsection{4. Conclusion}
THALAMUS provides the necessary structural foundation for stable, enterprise-scale reasoning. By integrating norma\-lization, isolation, and paral\-lelization, it ensures that the Knowledge Graph remains a coherent and high-integrity substrate for autonomous intel\-ligence. In CEREBRUM v2.52.0, THALAMUS has been extended with a hot-reload \texttt{/build} endpoint, a \texttt{ResearchAgent} feedback loop for human-approved edge ingestion, and a 1,357-test suite covering the full pipeline — confirming that high-throughput intelligent ingestion remains the stable foundation on which all reasoning capab\-ilities depend.

---

\subsection{5. Recent Advances (v2.52.0 -\textgreater  v2.52.0)}

THALAMUS has evolved from a prepr\-ocessing pipeline into a dynamic, bidir\-ectionally-connected ingestion layer since v2.51.1. The following describes key advances.

\textbf{Hot CSV Reload via /build Endpoint (Phase 54).} The \texttt{/build} endpoint enables runtime graph updates without server restart. A \texttt{POST /build} with a new CSV payload triggers THALAMUS to re-run the full ingestion pipeline — norma\-lization, dedup\-lication, embedding, structural encoding, and community detection — on the updated graph, then performs an atomic swap of the adapter's internal state. Active queries in flight during a \texttt{/build} operation are protected by query snapshot isolation (PAPER\_006, Phase 20) and complete against the pre-build graph state.

\textbf{ResearchAgent Feedback Loop (Phase 51).} The \texttt{ResearchAgent} is an autonomous agent that generates proposed KG triples by analyzing existing graph structure and querying external sources. Its proposals are surfaced to a human operator via a review queue. Upon approval, the approved triples are submitted to THALAMUS's \texttt{IngestionPipeline} as standard ingestion events — receiving full norma\-lization, dedup\-lication, namespace isolation, and confidence assignment. This closes the loop between autonomous reasoning (CORTEX) and structured knowledge ingestion (THALAMUS), enabling the graph to grow from its own reasoning activity.

\textbf{Full Pipeline Test Coverage.} The THALAMUS pipeline is now covered by 2177+ tests (up from 994 at v2.51.1). New test categories include:
\begin{itemize}
\item Streaming ingestion under high-velocity burst conditions
\item Namespace isolation regression tests (signal: vs text: collision prevention)
\item STDP discretizer integration tests within the pipeline
\item \texttt{/build} hot-reload atomicity and snapshot isolation tests
\item \texttt{ResearchAgent} approval-and-ingest workflow tests

\end{itemize}
\textbf{STDPDiscre\-tizer as Pipeline Stage (Phase 18).} The \texttt{STDPDiscretizer} is now an optional stage within \texttt{IngestionPipeline}, positioned between relation norma\-lization and the graph write-commit. Discretized causal edges emerge from the pipeline with \texttt{source="stdp"} provenance and a confidence score derived from the causal weight $w_{uv}$, making them first-class citizens of the graph's provenance model.

\textbf{Throughput Baseline Confirmed.} The 850\% throughput improvement (1,200 $\to$ 11,500 triples/sec) reported at v2.51.1 has been maintained through all pipeline additions. The hot-reload \texttt{/build} endpoint adds no steady-state latency to the ingestion path, as it operates on a separate execution context from the normal ingestion workers.

---
\subsection{Acknow\-ledgments}

The author gratefully ackno\-wledges the use of Claude (Anthropic) as a research assistant throughout this work. Claude assisted with mathe\-matical forma\-lization, code generation, manuscript preparation, and technical writing. All conceptual contr\-ibutions, archi\-tectural decisions, exper\-imental design, and intel\-lectual claims are solely the author's.

\textbf{References}
\begin{enumerate}
\item Doan, A., Halevy, A. Y., \& Ives, Z. G. (2012). Principles of Data Integration. Morgan Kaufmann.
\item Bizer, C., et al. (2009). Linked Data - The Story So Far. Intern\-ational Journal on Semantic Web and Information Systems.
\item Shvaiko, P., \& Euzenat, J. (2013). Ontology Matching: State of the Art and Future Challenges. IEEE TKDE.
\item Rahm, E., \& Bernstein, P. A. (2001). A survey of approaches to automatic schema matching. The VLDB Journal.
\item Getoor, L., \& Machan\-avajjhala, A. (2012). Entity resolution: Theory, practice \& open challenges. VLDB.
\item Christen, P. (2012). Data Matching: Concepts and Techniques for Record Linkage, Entity Resolution, and Duplicate Detection. Springer.
\item Buchorn, B. A. (2026). CEREBRUM v2.51: Complete Technical Specif\-ication for Autonomous Knowledge Graph Reasoning. [CEREBRUM\_REPORT\_PLACEHOLDER].

\end{enumerate}
---
\textbf{Reviewed on}: May 2, 2026 for version v2.52.0
